%% Introduction %%

Michel Melot établit le patrimoine comme une composante indispensable d'une communauté\footcite[]{melot2004}. En cela, un patrimoine ne nécessite, pour exister, que d'être désigné par une communauté comme représentatif de leur groupe. Or, dans ce mémoire, nous nous intéressons aux institutions patrimoniales, c'est-à-dire des organismes chargés de leur collecte, de leur conservation et de leur diffusion. Dans ce sens, la mise en place de ces institutions répond au besoin de ces communautés de désigner un gardien de leur patrimoine. Ces institutions peuvent être d'ordre privé comme l'est la \enquote{Société d’histoire et d’archéologie du Sedanais} qui prône la \enquote{Sauvegarde et mise en valeur du patrimoine historique et archéologique du pays sedanais}\footcite[]{sedan} ou encore comme le \enquote{Centre archives LGBTQI+} à Paris par lequel cette association LGBT+ exprime le besoin de \enquote{créer, collecter, classer, conserver et communiquer les archives LGBTQI+ à tous les publics}\footcite{lgbt}. Cela dit, nous resserrerons notre propos sur les institutions patrimoniales publiques, à laquelle appartient l'INA. Cette différenciation a une importance car si les institutions privées exercent librement leurs missions sans contrainte autre que celles qu'elles s'imposent, les institutions publiques sont sous l'autorité d'un État qui les désigne comme responsable d'un patrimoine national. En conséquent de quoi, leurs missions sont encadrées par des textes de loi désignant le patrimoine dont elles ont la responsabilité ainsi que le public auquel la communication du patrimoine conservé est destiné. C'est pourquoi il nous est essentiel de débuter cette analyse du contexte d'évolution du traitement documentaires des institutions patrimoniales en étudiant les lois qui les définissent et nous donnent une trace de leurs origines. C'est pour les respecter que les politiques documentaires appliquées répondent en premier aux prérogatives définit par ces textes de loi. 

\subsection{Avant l'INA, les cinémathécaires}

%% Post-INA %%

En 1949, dans l’ordonnance qui définit le statut de la RTF\footcite[]{zotero-item-335}, nous ne pouvons qu’interpréter dans ces mots une des futures missions de l’INA : 
\begin{displayquote}
	la Radiodiffusion-Télévision française [\dots] a seule qualité, dans les territoires de le République pour [\dots] radiodiffuser ses programmes et ou les mettre à la disposition d’autres organismes radiodiffuseur.
\end{displayquote}
En 1964, l’ORTF poursuit cette même mission\footcite[]{zotero-item-645} et 10 ans plus tard, quand Marceau Long présente le projet d’éclatement de l’ORTF, ces missions ont totalement disparues et l’INA ne figure pas encore dans les établissements publics résultant de la dissolution\footcite[][118]{cohen2021}. Ce n’est que dans la loi statuant de l’éclatement de l’ORTF, le 07 août 1974, que ses principales missions apparaissent enfin clairement et officiellement :  
\begin{displayquote}
	Il est créé un institut de l’audiovisuel chargé notamment de la conservation des archives, des recherches de création audiovisuelle et de la formation professionnelle.\footcite[]{zotero-item-331}
\end{displayquote}
Pourtant, dans cet intervalle, la mission de conservation et de communication des diffusions audiovisuelles n'est pas à l'abandon, elle est portée par celles que l’on appelait alors les « cinémathécaires ». Malgré le manque de reconnaissance social, politique et institutionnel de leurs activités\footcite[][182]{tible2023} elles posent les bases de que nous appelons aujourd’hui le traitement documentaire. Elles réalisent les « prises d’antenne », les indexations et recherchent les documents commandés par les chaînes de télévision de l'ORTF. Les informations sont sur fiches papier et sont indexées par thème ou par personnalité. On crée autant de fiches qu’il y a de thèmes indexés et elles sont rangées dans des meubles classeurs dédiés. Ces fiches doivent permettre à ces mêmes documentalistes de retrouver aisément les émissions demandées par les chaînes de télévision. En effet, la mission de communication est restreinte par la fragilité des supports audiovisuels. Il n'est pas envisageable de visionner plusieurs documents à chaque commande, autant du point de vue de la conservation que du temps consacrer à la recherche de contenu.

\subsection{Premières années d'existence : la commercialisation des images}
%% INA avant le dépôt légal %%

Ce sont donc sur ces bases que repose le traitement documentaire effectué à la création de l’INA. Cela dit, contrairement à ce que cela sous-entend quand on évoque la « conservation des archives », ces fiches ne sont pas empreintes de la volonté de patrimonialisation. L’INA vend ses documents et ses notices sont élaborées en fonction de ce besoin. Nous n'en sommes pas encore à désigner l'INA comme une \enquote{institution patrimoniale}. Cette vision des archives change progressivement au cours des années 70 et 80. La notion de patrimoine se popularise quand elle est employée par l'UNESCO en 1972 et la conservation des \enquote{images en mouvement} prend son importance avec les recommandations de l'UNESCO pour la sauvegarde et la conservation des \enquote{images en mouvement} le 27 octobre 1980\footcite[]{unesco}. En France c'est la loi dites « Léotard » en 1986\footcite[]{zotero-item-334} qui manifeste ce besoin d'introduire ces documents à la patrimonialisation. Elle ouvre la télédiffusion aux entreprises privées et appuie les missions de conservation et d’exploitation de l’INA :
\begin{displayquote}
	(INA), est chargé, [\dots]] de conserver et exploiter les archives audiovisuelles des sociétés nationales de programme. 
\end{displayquote}
Mais cette exploitation reste limitée aux chaînes de diffusions nationales et publiques, et uniquement à des fins professionnelles. 

\subsection{Le dépôt légal}
%% L'INA et le dépôt légal %%

C’est la loi qui implémente le dépôt légal de la radio et de la télévision, le 20 juin 1992\footcite[]{zotero-item-332} qui va durablement ajouter la patrimonialisation aux missions de l’INA et ouvrir l’exploitation des documents qu’elle conserve à des fins scientifiques\footcite[][23]{varra2023}.  
\begin{displayquote}
	L’institut national de l’audiovisuel est chargé de recueillir et de conserver les documents sonores et audiovisuels radiodiffusés ou télédiffusés, de participer à la diffusion des bibliographies nationales correspondantes et de mettre ces documents à la disposition du public pour consultation.
\end{displayquote}
C'est une étape importante dans l'histoire du traitement documentaire à l'INA. Le traitement de ces documents doit désormais s’adapter à une exploitation commerciale et scientifique. D’autant plus que l’INA ne peut vendre que des programmes sur lesquels elle a les droits, ce qui n’inclue pas tous les documents conservés au titre du dépôt légal. Ces documents constituent une exception au droit d’auteur. L’INA a le droit de les exploiter et de les diffuser uniquement à des fins scientifiques. Cette multiplication des missions donne naissance à une modification importante du fonctionnement de l’INA\footcite[][23]{varra2023}. L'usage des documents que conservent l'INA change radicalement et la description doit s'y adapter. Avec le dépôt légal, il n'est plus seulement question de créer des fiches pour les documents permettant de retrouver des thèmes ou des personnalités. En conséquence, est créé à l’INA une nouvelle direction pour héberger cette nouvelle mission patrimoniale et dès 1995 est ouvert le centre de consultation du dépôt légal à la Bibliothèque François Mitterrand. Cette loi, marquant un changement radical dans le traitement documentaire, est le parfait exemple d'une loi impactant le public cible du traitement documentaire et en conséquent ce qu'il décrit dans le document. Par la suite, le champ du dépôt sera précisé par un décret\footcite[]{zotero-item-330} et continuera de s’étendre à d’autres chaînes et d’autres programmes\footcite[][25]{varra2023}, notamment avec l’arrivée de la TNT, puis des programmes télévisés sur le web\footcite[]{zotero-item-647}. En 2004, le dépôt légal est intégré au code du patrimoine\footcite[]{2004}, qui remplace la loi de 1992.

%% Le dépôt légal dans d'autres institutions %%

Il y a encore quelques années, ces documents audiovisuels étaient uniquement conservés, aux yeux de la loi, dans un objectif de consommation. Avec le dépôt légal, ces documents audiovisuels font l’objet d’une nouvelle injonction à la patrimonialisation. Ils s’inscrivent désormais dans un patrimoine national\footcite[][74]{bermès2024} et les politiques documentaires de l’INA s’en voient questionnées. Là où le public cible était unique, il est désormais multiple et les impératifs métiers concernant le traitement documentaire doivent s’adapter à des prérogatives légales. En effet, il est dans la nature du dépôt légal, dès sa création en 1537, de léguer un discours national aux générations futures. En s'appliquant aux documents audiovisuels télévisés, le dépôt légal l'inscrit dans le patrimoine national français. D’abord destiné aux livres, le dépôt s’étendra au fur et à mesure à plusieurs supports permettant la transmission de la pensée\footcite[][18-19]{saby2013}. Aujourd’hui, les institutions patrimoniales qui sont chargées du dépôt légal aux côtés de l’INA sont le CNC pour l’image animée\footcite[]{zotero-item-673} et la BnF pour les livres périodiques, documents cartographiques, iconographiques ainsi que la musique imprimée et les documents sonores et vidéogrammes\footcite[]{zotero-item-675}. Pourtant, malgré cette injonction légale commune, ces institutions ne proposent pas le même traitement documentaire, elles répondent individuellement aux restrictions qu’imposent leurs supports mais elles s’appuient également sur leur histoire. Toutes ces institutions n’ont ni crée le traitement documentaire, ni attendu le dépôt légal pour établir des méthodes de traitement, comme nous l’avons vu avec les cinémathécaires de l’ORTF. La principale différence réside dans l’héritage dont ils disposent. En 1537, les livres n’étaient pas des supports récents, au contraire des diffusions audiovisuelles, télévisées et radiophoniques. Ces supports n’ont pas hérité de centaines d’années de réflexion sur leur traitement documentaire et leur catalogage. Bien qu'elles puissent s'inspirer des traitements appliqués à d'autres supports. Le dépôt légal n’a pas fait qu’institutionnaliser leur conservation, il a institutionnalisé la réflexion sur leur traitement dans un objectif de communication scientifique et de patrimonialisation. 

Dans ce cadre, nous pouvons questionner le traitement documentaire dans un cadre où le dépôt légal n’a pas institutionnalisé la conservation de ses diffusions télévisées. En Suisse, malgré une loi nommée « dépôt légal », la RTS ne conserve que des émissions\footcite[][11]{griveau2024}, et non l’entièreté du flux, 24 heures sur 24, comme le fait l’INA depuis 2001 grâce à un dispositif de captation numérique par liaisons satellites et fibre optique\footcite[][23]{saby2013}. Le flux télévisuel n’est donc pas conservé. Les émissions de leur fond ont été sélectionnées en fonction de critères documentaires tels que leur réutilisation\footcite[][16]{griveau2024}, conservant une logique commerciale. Actuellement, les fonds de la RTS sont de 680 000 heures\footcite[][13]{griveau2024} contre 28 172 492 heures au titre du dépôt légal le 31 décembre 2025\footcite[]{zotero-item-655} à l'INA. Aux États-Unis la librairie du congrès ne conserve que 3.5 millions d’heures d’enregistrements de toutes origines, y compris la radio\footcite[]{zotero-item-654} et une partie des archives audiovisuelles des « National Archives » manquent encore de description\footnote{(Voir entretien de Pascal Lebeurier sur l'expertise du documentaliste~\ref{ann:transcriptions},\nameref{ann:transcriptions}), p.~\pageref{ann:transcriptions}.}.

%% Conclusion et ouverture à la sous-partie suivante %%

L'exemple de l'INA est une bonne représentation de l'impact que peut avoir le cadre légal sur une institution. À la naissance de la conservation des supports audiovisuels télévisés par les cinémathécaires de l'ORTF, la notion conservation patrimoniale est de moindre importance en comparaison de l'importance portée à la conservation à but commerciale. Les informations que l'on retrouve dans les fiches ont pour objectif de répondre à ces usages. Elles doivent permettre de trouver rapidement des images pour illustrer les émissions. Ce n'est que quand le dépôt légal est appliqué à la télévision que l'institut national de l’audiovisuel est considéré comme un institut véritablement à but patrimonial. Le traitement documentaire est désormais destiné à un nouvel usage scientifique où la rapidité à trouver une information importe moins que de transmettre une description de l'ensemble du flux télévisé afin de répondre à de nouveaux usages, de nouveaux besoins, de nouvelles pratiques, sans toujours savoir ce que ce nouveau public attendra. Ce nouveau service du dépôt légal doit associer les mêmes contraintes de fragilité des supports avec un nouveau public pour lequel la consultation des documents s'inscrit dans leurs pratiques. D'autant plus que cela s'accompagne d'une massification des collections de l'INA et des données les décrivant. Nous décrirons plus en détail cette massification des données et comment ils impactent le traitement documentaire. Mais avant cela, pour appuyer sur l'impact majeur qu'a cette loi du dépôt légal sur le traitement documentaire, il est important d'appuyer les changements de statut des documents audiovisuels que cela entraîne.